\chapter{Operazioni sui Linguaggi regolari}

Per descrivere completamente la classe dei linguaggi regolari, oltre alla definizione di automa, è necessario introdurre alcune operazioni che permettono di costruire nuovi linguaggi regolari a partire da altri.

In particolare vedremo come i linguaggi regolari siano chiusi rispetto alle operazioni di unione, concatenazione e iterazione.

Per dimostrare ciò più agilmente però, ci sarà necessaria una nuova classe di automi, equivalente ai DFA, ovvero i \textbf{DFA non-restarting}.

\section{DFA Non-Restarting}

\dfn{DFA Non-restarting}{
  Un DFA \(\mcM\) è \textbf{Non-Restarting} se \(\mcM=(Q,A,\delta,q_1,F)\) con \(q_1 \nin \delta(QxA)\), ovvero non è possibile tornare allo stato iniziale dopo aver letto un simbolo di input.
}

Una definizione analoga può essere data per gli NFA, sincerandosi che \(q_1 \nin \bigcup\limits_{Q' \in \delta(Q \times A)} Q'\).

Come anticipato, vediamo ora come questa classe di automi sia equivalente a quella dei DFA

\begin{halfframedbox}{red!75!black}{DFA Non-Restarting}
  \textbf{Enunciato. } \(L\subseteq A\) regolare \(\implies\) % chktex 21
   esiste un DFA \textbf{non-restarting} \(\mcM'\) tale che \(L=L(\mcM')\).

  \begin{center}
    \textcolor{red!75!black}{\hdashrule[0.5ex]{18cm}{1pt}{3mm 3pt 1mm 2pt}}
  \end{center}

  \textbf{Dimostrazione. } Sia \(\mcM=(Q,A,\delta,q_1,F)\) un DFA tale che \(L=L(\mcM)\) e sia \(q_0 \nin Q\).

  Definiamo \(\mcM'=(Q\cup\cbrakets{q_0},A,\delta',q_0,F')\) dove \(\delta'\) e \(F'\) sono rispettivamente:

  \begin{equation}
    \delta'(q,a)=\begin{cases}
      \delta(q,a), & q \in Q \land a \in A\\
      \delta(q_1,a), &\text{altrimenti}
    \end{cases}
  \end{equation}

  Che è quindi identica a \(\delta\) tranne per \(q_0\) che emula l'uscita da \(q_1\).

  \begin{equation}
    F'=\begin{cases}
      F, & q_1 \nin F\\
      F \cup \cbrakets{q_0}, & q_1 \in F
    \end{cases}
  \end{equation}

  Evidentemente \(\mcM'\) è un DFA \textbf{non-restarting} e \(L(\mcM')=L(\mcM)\).
\end{halfframedbox}

\section{Chiusura dei Linguaggi Regolari}

\subsection{Unione}

Grazie a questa nuova classe di automi possiamo dimostrare le prime proprietà di chiusura dei linguaggi regolari. Iniziamo vedendo come la classe dei linguaggi regolari sia chiusa rispetto all'operazione di \textbf{unione}.

\begin{halfframedbox}{red!75!black}{Regolarità dell'Unione}\label{Regolarità dell'unione}
  \textbf{Enunciato. } Siano \(L,L' \subseteq A^*\) regolari. Allora \(L \cup L'\) è regolare.

  \begin{center}
    \textcolor{red!75!black}{\hdashrule[0.5ex]{18cm}{1pt}{3mm 3pt 1mm 2pt}}
  \end{center}

  \textbf{Dimostrazione. } Siano \(\mcM=(Q,A,\delta,q_1,F)\) e \(\mcM'=(Q',A,\delta',q_1',F')\) DFA \textbf{non-restarting} tali che \(L=L(\mcM)\) e \(L'=L(\mcM')\) con \(Q \cap Q' = \varnothing\).

  Definiamo un NFA \(\widehat{\mcM}=(Q\cup Q'\setminus\cbrakets{q_1,q_1'}\cup\cbrakets{q_0},A,\widehat{\delta},q_0,\widehat{F})\) con
  \begin{equation}
    \begin{split}
      \widehat{F}&=\begin{cases}
        F\cup F', & q_1 \nin F \land q_1' \nin F'\\
        F\cup F'\setminus\cbrakets{q_1,q_1'}\cup\cbrakets{q_0}, & q_1 \in F \lor q_1' \in F'\\
      \end{cases} \\
      \forall a \in A: \widehat{\delta}(q,a)&=\begin{cases}
        \cbrakets{\delta(q,a)}, & q \in Q \\
        \cbrakets{\delta'(q,a)}, & q \in Q' \\
        \cbrakets{\delta(q_1,a), \delta'(q_1',a)}, & q=q_0
      \end{cases}
    \end{split}
  \end{equation}

  È evidente che il linguaggio accettato da questo NFA è \(L(\widehat{\mcM})=L(\mcM)\cup L(\mcM')=L\cup L'\).
\end{halfframedbox}

È possibile ora facilmente vedere come i linguaggi regolari siano chiusi anche rispetto all'intersezione.

\begin{halfframedbox}{red!75!black}{Corollario: Regolarità dell'Intersezione}\label{Regolarità dell'intersezione}
  \textbf{Enunciato. } Siano \(L,L' \subseteq A^*\) regolari. Allora \(L \cap L'\) è regolare.

  \begin{center}
    \textcolor{red!75!black}{\hdashrule[0.5ex]{18cm}{1pt}{3mm 3pt 1mm 2pt}}
  \end{center}

  \textbf{Dimostrazione. } Si ha \(L\cap L'=\overline{\rbrakets{\overline{L}\cup\overline{L'}}}\). Dai teoremi~\ref{Regolarità del Complemento} e~\ref{Regolarità dell'unione} abbiamo la tesi.

\end{halfframedbox}

Vediamo ora un applicazione pratica di questi due teoremi.

Si consideri \(L_1=\cbrakets{w \in \cbrakets{a,b}* \vert w \text{ contiene almeno 2 a}}\) e \(L_1=\cbrakets{w \in \cbrakets{a,b}* \vert w \text{ contiene almeno 2 b}}\). Si consideri gli automi ottenuti nell'esercizio~\ref{Esercizio NFA 2a2b}.

Avremo che \(\overline{L_1}=L(\mcM_1)\) e \(\overline{L_2}=L(\mcM_2)\) dove \(\mcM_1\) e \(\mcM_2\) sono rispettivamente:

\begin{figure}[ht]
  \centering
  \subfloat[Diagramma di \(\mcM_1\)]{
    \begin{tikzpicture}[initial text=, every state/.style={draw=blue!50,very thick,fill=blue!20}]
      \node [state, initial, accepting]  (q_1) {\(q_1\)};
      \node [state, right = of q_1, accepting] (q_2) {\(q_2\)};
      \node [state, below = of q_2, accepting] (q_3) {\(q_3\)};
      \node [state, left = of q_3] (q_4) {\(q_4\)};

      \path [-latex, every loop/.style = {-latex}]
      (q_1) edge [right] node [above] {b} (q_2)
      (q_1) edge [right] node [above] {a} (q_3)
      (q_2) edge [right] node [right] {a} (q_3)
      (q_2) edge [loop right] node [right] {b} (q_2)
      (q_3) edge [left] node [above] {a} (q_4)
      (q_3) edge [loop right] node [right] {b} (q_3)
      (q_4) edge [loop left] node [left] {a,b} (q_4);
    \end{tikzpicture}
  }
  \subfloat[Diagramma di \(\mcM_2\)]{
    \begin{tikzpicture}[initial text=, every state/.style={draw=blue!50,very thick,fill=blue!20}]
      \node [state, initial, accepting]  (q_1) {\(q_1\)};
      \node [state, right = of q_1, accepting] (q_2) {\(q_2\)};
      \node [state, below = of q_2, accepting] (q_3) {\(q_3\)};
      \node [state, left = of q_3] (q_4) {\(q_4\)};

      \path [-latex, every loop/.style = {-latex}]
      (q_1) edge [right] node [above] {a} (q_2)
      (q_1) edge [right] node [above] {b} (q_3)
      (q_2) edge [right] node [right] {b} (q_3)
      (q_2) edge [loop right] node [right] {a} (q_2)
      (q_3) edge [right] node [above] {b} (q_4)
      (q_3) edge [loop right] node [right] {a} (q_3)
      (q_4) edge [loop left] node [left] {a,b} (q_4);
    \end{tikzpicture}
  }
\end{figure}

Avremo che \(\overline{L_1}\cup\overline{L_2}=L(\widehat{\mcM})\) dove \(\widehat{\mcM}\) sarà:

\begin{figure}[ht]
  \centering
  \subfloat[Diagramma di \(\widehat{\mcM}\)]{
    \begin{tikzpicture}[initial text=, every state/.style={draw=blue!50,very thick,fill=blue!20}]
      \node [state, initial, accepting]  (q_0) {\(q_0\)};
      \node [state, above right = of q_0, accepting] (q_1) {\(q'_1\)};
      \node [state, below right = of q_0, accepting] (q_2) {\(q_1\)};
      \node [state, right = of q_1, accepting] (q_3) {\(q_2'\)};
      \node [state, right = of q_3] (q_4) {\(q_3'\)};
      \node [state, right = of q_2, accepting] (q_5) {\(q_2\)};
      \node [state, right = of q_5] (q_6) {\(q_3\)};

      \path [-latex, every loop/.style = {-latex}]
      (q_0) edge [above] node [left] {a} (q_1)
      (q_0) edge [below] node [left] {b} (q_2)
      (q_0) edge [right] node [below] {b} (q_3)
      (q_0) edge [above] node [above] {a} (q_5)
      (q_1) edge [right] node [above] {b} (q_3)
      (q_1) edge [loop above] node [above] {a} (q_1)
      (q_2) edge [loop below] node [below] {b} (q_2)
      (q_2) edge [right] node [above] {a} (q_5)
      (q_3) edge [loop above] node [above] {a} (q_3)
      (q_3) edge [right] node [above] {b} (q_4)
      (q_4) edge [loop above] node [above] {a, b} (q_4)
      (q_5) edge [loop below] node [below] {b} (q_5)
      (q_5) edge [right] node [above] {a} (q_6)
      (q_6) edge [loop below] node [below] {a, b} (q_6);
    \end{tikzpicture}
  }
  \subfloat[Diagramma di \(\widehat{\mcM}\) semplificato]{
    \begin{tikzpicture}[initial text=, every state/.style={draw=blue!50,very thick,fill=blue!20}]
      \node [state, initial, accepting]  (q_0) {\(q_0\)};
      \node [state, above right = of q_0, accepting] (q_1) {\(q'_1\)};
      \node [state, below right = of q_0, accepting] (q_2) {\(q_1\)};
      \node [state, right = of q_1, accepting] (q_3) {\(q_2'\)};
      \node [state, right = 2.65 of q_0] (q_4) {\(q_3\)};
      \node [state, right = of q_2, accepting] (q_5) {\(q_2\)};

      \path [-latex, every loop/.style = {-latex}]
      (q_0) edge [above] node [left] {a} (q_1)
      (q_0) edge [below] node [left] {b} (q_2)
      (q_0) edge [right] node [below] {b} (q_3)
      (q_0) edge [above] node [above] {a} (q_5)
      (q_1) edge [right] node [above] {b} (q_3)
      (q_1) edge [loop above] node [above] {a} (q_1)
      (q_2) edge [loop below] node [below] {b} (q_2)
      (q_2) edge [right] node [above] {a} (q_5)
      (q_3) edge [loop above] node [above] {a} (q_3)
      (q_3) edge [below] node [right] {b} (q_4)
      (q_4) edge [loop right] node [right] {a, b} (q_4)
      (q_5) edge [loop below] node [below] {b} (q_5)
      (q_5) edge [above] node [right] {a} (q_4);
    \end{tikzpicture}
  }
\end{figure}

A questo punto usando la tecnica usata nel teorema~\ref{Corrispondenza tra DFA e NFA} possiamo ottenere un DFA da complementare.

\newpage

Andiamo ora a elencare alcuni risultati derivanti dalla chiusura dei linguaggi regolari rispetto alle operazioni appena viste.

\begin{halfframedbox}{red!75!black}{Proposizione: Regolarità dell'insieme vuoto}
  \textbf{Enunciato. } \(\varnothing \subseteq A^*\) è regolare.

  \begin{center}
    \textcolor{red!75!black}{\hdashrule[0.5ex]{18cm}{1pt}{3mm 3pt 1mm 2pt}}
  \end{center}

  \textbf{Dimostrazione. } Accettato da qualsiasi DFA con \(F=\varnothing\).
\end{halfframedbox}

\begin{halfframedbox}{red!75!black}{Proposizione: Regolarità dei singleton}
  \textbf{Enunciato. } \(\forall w \in A^*, \cbrakets{w}\) è regolare.

  \begin{center}
    \textcolor{red!75!black}{\hdashrule[0.5ex]{18cm}{1pt}{3mm 3pt 1mm 2pt}}
  \end{center}

  \textbf{Dimostrazione. } Se \(w=w_1\cdots w_n\), con \(w_1 \in A\) e \(n \in \mbN\), \(L=\cbrakets{w}\) sarà accettato dall'NFA con diagramma:

  \centering
  \begin{tikzpicture}[initial text=, every state/.style={draw=blue!50,very thick,fill=blue!20}]
    \tikzset{dotsstyle/.style={draw=none,fill=none}}

    \node [state, initial] (q_0) {\(q_0\)};
    \node [state, right = of q_0] (q_1) {\(q_1\)};
    \node [state, dotsstyle, right = of q_1] (q_2) {\(\cdots\)};
    \node [state, right = of q_2, accepting] (q_3) {\(q_i\)};

    \path [-latex, every loop/.style = {-latex}]
    (q_0) edge [right] node [above] {\(w_1\)} (q_1)
    (q_1) edge [right] node [above] {\(w_2\)} (q_2)
    (q_2) edge [right] node [above] {\(w_n\)} (q_3);
  \end{tikzpicture}
\end{halfframedbox}

\begin{halfframedbox}{red!75!black}{Corollario: Regolarità dei linguaggi finiti}\label{Regolarità dei linguaggi finiti}
  \textbf{Enunciato. } \(L\) finito \(\implies L\) regolare.

  \begin{center}
    \textcolor{red!75!black}{\hdashrule[0.5ex]{18cm}{1pt}{3mm 3pt 1mm 2pt}}
  \end{center}

  \textbf{Dimostrazione. } Tutti i linguaggi finiti sono unione finita di singleton. Dalla proposizione appena vista e dalla chiusura per unione abbiamo la tesi.
\end{halfframedbox}

\subsection{Prodotto per Concatenazione}

\dfn{Prodotto per Concatenazione}{
  Se \(L,L'\subseteq A^*\), definiamo:
  \begin{equation}
    LL'=\cbrakets{ww' \in A^* \vert w \in L \land w' \in L'}
  \end{equation}
}

È importante notare che questo linguaggio conterrà parole di uno dei due linguaggi originali \textit{se e solo se} l'altro contiene \(\varepsilon\).

\begin{halfframedbox}{red!75!black}{Chiusura dei Linguaggi regolari per Prodotto}
  \textbf{Enunciato. } Siano \(L,L' \subseteq A^*\) regolari. Allora \(LL'\) è regolare.

  \begin{center}
    \textcolor{red!75!black}{\hdashrule[0.5ex]{18cm}{1pt}{3mm 3pt 1mm 2pt}}
  \end{center}

  \textbf{Dimostrazione. } Siano \(\mcM=(Q,A,\delta,q_1,F)\) e \(\mcM'=(Q',A,\delta',q_1',F')\) DFA \textbf{non-restarting} tali che \(L=L(\mcM)\) e \(L'=L(\mcM')\) con \(Q \cap Q' = \varnothing\).

  Definiamo un NFA \(\widehat{\mcM}=(Q\cup Q', A, \widehat{\delta}, q_1, \widehat{F})\) con

  \begin{equation}
    \begin{split}
      \widehat{F}&=\begin{cases}
        F\cup F', & q_1' \in F'\\
        F', & q_1' \nin F'\\
      \end{cases} \\
      \forall a \in A: \widehat{\delta}(q,a)&=\begin{cases}
        \cbrakets{\delta(q,a)}, & q \in Q \setminus F\\
        \cbrakets{\delta'(q,a)}, & q \in Q' \\
        \cbrakets{\delta(q,a), \delta'(q_1',a)}, & q\in F
      \end{cases}
    \end{split}
  \end{equation}

  È evidente che il linguaggio accettato da questo NFA è \(L(\widehat{\mcM})=L(\mcM)L(\mcM')=LL'\).

\end{halfframedbox}

\newpage

\begin{generalbox}[colframe=azure-gradient-3!90!black]{Esempio: Concatenazione di Linguaggi Regolari}

  Consideriamo ad esempio il linguaggio \(L_1L_1\) con \(L_1\) definito come in precedenza. Consideriamo due automi \(\mcM\) e \(\mcM'\) che accettano \(L_1\), i quali avranno entrambi la forma:

  \begin{center}
      \begin{tikzpicture}[initial text=, every state/.style={draw=blue!50,very thick,fill=blue!20}]
        \node [state, initial]  (q_1) {\(q_1\)};
        \node [state, right = of q_1] (q_2) {\(q_2\)};
        \node [state, right = of q_2, accepting] (q_3) {\(q_3\)};

        \path[-latex, every loop/.style = {-latex}]
        (q_1) edge [loop above]  node [above] {b} (q_1)
        (q_1) edge [above]  node [above] {a} (q_2)
        (q_2) edge [above]  node [above] {a} (q_3)
        (q_2) edge [loop above]  node [above] {b} (q_1)
        (q_3) edge [loop above]  node [above] {a,b} (q_3);
      \end{tikzpicture}
  \end{center}

  Si avrà che il prodotto per concatenazione \(L_1L_1\) sarà accettato dall'automa \(\widehat{\mcM}\):

  \begin{center}
    \begin{tikzpicture}[initial text=, every state/.style={draw=blue!50,very thick,fill=blue!20}]
      \node [state, initial]  (q_1) {\(q_1\)};
      \node [state, right = of q_1] (q_2) {\(q_2\)};
      \node [state, right = of q_2] (q_3) {\(q_3\)};
      \node [state, below = of q_3] (q_4) {\(q_1'\)};
      \node [state, right = of q_4] (q_5) {\(q_2'\)};
      \node [state, right = of q_5, accepting] (q_6) {\(q_3'\)};

      \path[-latex, every loop/.style = {-latex}]
      (q_1) edge [loop above]  node [above] {b} (q_1)
      (q_1) edge [above]  node [above] {a} (q_2)
      (q_2) edge [above]  node [above] {a} (q_3)
      (q_2) edge [loop above]  node [above] {b} (q_1)
      (q_3) edge [loop above]  node [above] {a,b} (q_3)
      (q_3) edge [below] node [right] {b} (q_4)
      (q_3) edge [above] node [right] {a} (q_5)
      (q_4) edge [loop below]  node [below] {a} (q_4)
      (q_4) edge [right]  node [below] {b} (q_5)
      (q_5) edge [loop below]  node [below] {b} (q_5)
      (q_5) edge [right]  node [below] {a} (q_6)
      (q_6) edge [loop below]  node [below] {a,b} (q_6);
    \end{tikzpicture}
  \end{center}
\end{generalbox}


In particolare questo prodotto è un esempio di prodotto per concatenazione di un linguaggio per se stesso. Possiamo in generale definire dal prodotto l'operazione potenza di un linguaggio come:

\begin{equation}
  \begin{cases}
    L^0=\cbrakets{\varepsilon}\\
    L^{n+1}=LL^n
  \end{cases}
\end{equation}

\subsection{Iterazione}

\dfn{Iterazione (Star)}{
  Se \(L \subseteq A^*\), definiamo:
  \begin{equation}
    L^*=\cbrakets{w_1\cdots w_n \vert n \in \mbN, w_i \in L}=\bigcup\limits_{n \in \mbN}^\infty  L^n
  \end{equation}
}

Essendo l'operazione Star definita tramite uso di un unione infinita di linguaggi regolari, non possiamo ricavarla direttamente come applicazione finita di operazioni di unione e prodotto. Per questo è necessario definire separatamente questa operazione.

\begin{halfframedbox}{red!75!black}{Chiusura dei Linguaggi regolari per Iterazione}\label{Regolarità dell'iterazione}
  \textbf{Enunciato. } Sia \(L \subseteq A^*\) regolare. Allora \(L^*\) è regolare.

  \begin{center}
    \textcolor{red!75!black}{\hdashrule[0.5ex]{18cm}{1pt}{3mm 3pt 1mm 2pt}}
  \end{center}

  \textbf{Dimostrazione. } Sia \(\mcM=(Q,A,\delta,q_1,F)\) un DFA \textbf{non-restarting} tale che \(L=L(\mcM)\).

  Definiamo un NFA \(\widehat{\mcM}=(Q, A, \widehat{\delta}, q_1, \cbrakets{q_1})\) con:

  \begin{equation}
      \forall a \in A: \widehat{\delta}(q,a)=\begin{cases}
        \cbrakets{\delta(q,a)}, & \delta(q,a) \nin F\\
        \cbrakets{\delta(q,a), q_1}, \text{ altrimenti}
      \end{cases}
  \end{equation}

  Per vedere come \(L^*=L(\mcM)\) verifichiamo la doppia inclusione.

  \begin{itemize}
    \item [``\(\subseteq\)''] Verifichiamo per induzione che \(L^*\subseteq L(\widehat{\mcM})\). Come base induttiva possiamo facilmente notare che \(\varepsilon \in L(\widehat{\mcM})\).
    Per quanto riguarda il passo induttivo invece supponiamo \(L^{n-1}\subseteq L(\widehat{\mcM})\) e mostriamo che \(L^n \subseteq L(\widehat{\mcM})\).

    Siano \(u \in L^{n-1}, v \in L\) tali che \(w=uv\). Si ha \(q_1 \in \widehat{\delta}^*(q_1,u)\) per ipotesi induttiva. Ma allora essendo \(v \in L\) partendo da \(q_1\) si tornerà in \(q_1\) per costruzione di \(\widehat{\mcM}\). Si avrà dunque che \(q_1 \in \widehat{\delta}^*(q_1,uv) \implies L^n \subseteq L(\widehat{\mcM})\)

    \item [``\(\supseteq\)''] Si ha che \(L(\widehat{\mcM}) \subseteq L^*\) poiché, per costruzione di \(\widehat{\mcM}, w \in L(\mcM)\) è composizione di parole di \(L\), dovendo per terminare raggiungere \(q_1\), ovvero l'unico stato terminale di \(\mcM\). Essendo che \(L\) era \textbf{non-restarting} l'unico modo per raggiungere \(q_1\) in \(\widehat{\mcM}\) è raggiungere uno stato terminale di \(L\).
  \end{itemize}

\end{halfframedbox}

\newpage

\begin{generalbox}
  [colframe=azure-gradient-3!90!black]
  {Esempio: Iterazione di Linguaggi Regolari}
  Consideriamo ad esempio il linguaggio \(L=\cbrakets{ab,ba}\), accettato dall'automa \(\mcM\):

  \medskip
  \begin{center}
    \begin{tikzpicture}[initial text=, every state/.style={draw=blue!50,very thick,fill=blue!20}]
      \node [state, initial]  (q_1) {\(q_1\)};
      \node [state, above right = of q_1] (q_2) {\(q_2\)};
      \node [state, below right = of q_1] (q_3) {\(q_3\)};
      \node [state, above right = of q_3, accepting] (q_4) {\(q_4\)};
      \node [state, below = 0.7 of q_2] (q_5) {\(q_5\)};

      \path[-latex, every loop/.style = {-latex}]
      (q_1) edge [above]  node [above] {a} (q_2)
      (q_1) edge [below]  node [below] {b} (q_3)
      (q_2) edge [above]  node [above] {b} (q_4)
      (q_2) edge [below] node [right] {a} (q_5)
      (q_3) edge [below]  node [below] {a} (q_4)
      (q_3) edge [above] node [right] {b} (q_5)
      (q_4) edge [left] node [above] {a,b} (q_5)
      (q_5) edge [loop left] node [below right] {a,b} (q_5);
    \end{tikzpicture}
  \end{center}

  Si avrà che l'iterazione \(L^*\) sarà accettata dall'automa \(\widehat{\mcM}\):

  \medskip
  \begin{center}
    \begin{tikzpicture}[initial text=, every state/.style={draw=blue!50,very thick,fill=blue!20}]
      \node [state, initial, accepting]  (q_1) {\(q_1\)};
      \node [state, above right = of q_1] (q_2) {\(q_2\)};
      \node [state, below right = of q_1] (q_3) {\(q_3\)};
      \node [state, above right = of q_3] (q_4) {\(q_4\)};
      \node [state, below = 0.7 of q_2] (q_5) {\(q_5\)};

      \path[-latex, every loop/.style = {-latex}]
      (q_1) edge [above]  node [above] {a} (q_2)
      (q_1) edge [below]  node [below] {b} (q_3)
      (q_2) edge [above]  node [above] {b} (q_4)
      (q_2) edge [below] node [right] {a} (q_5)
      (q_2) edge [bend right, above] node [above] {b} (q_1)
      (q_3) edge [below]  node [below] {a} (q_4)
      (q_3) edge [above] node [right] {b} (q_5)
      (q_3) edge [bend left, below] node [below] {a} (q_1)
      (q_4) edge [left] node [above] {a,b} (q_5)
      (q_5) edge [loop left] node [below right] {a,b} (q_5);
    \end{tikzpicture}
  \end{center}

  È evidente però che questo automa non sia il più semplice possibile che accetti \(L^*\). Infatti, \(\widehat{\mcM}\) può essere semplificato in:

  \medskip
  \begin{center}
    \begin{tikzpicture}[initial text=, every state/.style={draw=blue!50,very thick,fill=blue!20}]
      \node [state, initial, accepting]  (q_1) {\(q_1\)};
      \node [state, above right = of q_1] (q_2) {\(q_2\)};
      \node [state, below right = of q_1] (q_3) {\(q_3\)};

      \path[-latex, every loop/.style = {-latex}]
      (q_1) edge [above]  node [above] {a} (q_2)
      (q_1) edge [below]  node [below] {b} (q_3)
      (q_2) edge [bend right, above] node [above] {b} (q_1)
      (q_3) edge [bend left, below] node [below] {a} (q_1);
    \end{tikzpicture}
  \end{center}
\end{generalbox}

\newpage

\begin{halfframedbox}{red!75!black}{Teorema di Kleene}\label{Teorema di Kleene}

  \textbf{Enunciato.} Un linguaggio \(L \subseteq A^*\) è regolare \textit{se e solo se} è finito oppure è ottenuto da linguaggi finiti mediante un numero finito di operazioni di unione, concatenazione e iterazione.

  \begin{center}
    \textcolor{red!75!black}{\hdashrule[0.5ex]{18cm}{1pt}{3mm 3pt 1mm 2pt}}
  \end{center}

  \textbf{Dimostrazione.} Dimostriamo la doppia implicazione.

  \begin{itemize}
    \item [``\(\Rightarrow\)''] Ovvio dai teoremi~\ref{Regolarità dei linguaggi finiti}~\ref{Regolarità dell'unione},~\ref{Regolarità dell'intersezione} e~\ref{Regolarità dell'iterazione}.
    \item [``\(\Leftarrow\)''] Sia \(\mcM=(Q,A,\delta,q_1,F)\) un DFA con \(Q=\cbrakets{q_1,\cdots,q_n}\). Definiamo dunque una famiglia di linguaggi della forma \(R^{(k)}_{i,j}\), dove:
    \begin{equation}
      \begin{split}
        \forall 1 \leq i&, j \leq n \land 0 \leq k \leq n:\\
        &R^{(k)}_{i,j}=\cbrakets{w \in A^* \vert \delta^*(q_i,w)=q_j \land \delta^*(q_i,w') \in \cbrakets{q_1,\cdots,q_k} \forall w'\neq \varepsilon \text{ prefisso proprio di } w }
      \end{split}
    \end{equation}

    Ovvero fissati \(i\), \(j\) e \(k\), il linguaggio \(R^{(k)}_{i,j}\) conterrà le parole che mi permettono di raggiungere lo stato \(q_j\) da \(q_i\) passando solo per stati \(q \in \cbrakets{q_1,\cdots,q_k}\).

    Nella definizione formale data, la seconda condizione richiede sostanzialmente che ogni sotto stringa di \(w\) ci porti \textbf{esclusivamente} in stati da \(q_1\) a \(q_k\).

    È per questo fondamentale specificare che \(w'\) debba essere un prefisso \textbf{proprio} di \(w\), ovvero diverso da \(w\) stesso, poiché altrimenti ogni insieme con \(k<j\) sarebbe vuoto, poiché le parole al suo intero dovrebbero raggiungere lo stato \(q_j\) dovendo passare esse stesse, non solo le loro sotto stringhe, solo per stati da \(q_1\) a \(q_k\), ovvero impossibilitando l'arrivo in \(q_j\).

    Analogamente è fondamentale specificare che \(w'\) non sia la parola vuota, per evitare una contraddizione simile a quella appena vista ma nel caso \(i>k\). Si avrà in fatti in questo caso che per ogni parola, si dovrà avere che:

    \begin{equation}
      \cbrakets{q_1, \cdots, q_k} \ni \delta^*(q_i,\varepsilon) = q_i \notin \cbrakets{q_1, \cdots, q_k}
    \end{equation}

    Definiti dunque questi insiemi possiamo osservare che:

    \begin{equation}
      \begin{split}
        R^{(0)}_{i,i}&=\cbrakets{\varepsilon} \cup \cbrakets{a \in A^* \vert \delta(q_i, a) = q_i} \text{ è un insieme finito}\\
        \text{Con } i \neq j: R^{(0)}_{i,j}&=\cbrakets{a \in A \vert \delta(q_i, a) = q_j}, \text{ è un insieme finito} \\
        R^{(k+1)}_{i,j}&=R^{(k)}_{i,j} \cup R^{(k)}_{i,k+1}\cdot {(R^{(k)}_{k+1,k+1})}^* \cdot R^{(k)}_{k+1,j}
      \end{split}
    \end{equation}


  In particolare l'ultima osservazione ci dice che tutte le parole che ci portano da \(q_i\) a \(q_j\) attraversando esclusivamente stati da \(q_1\) a \(q_{k+1}\) possono essere di due tipi, ovvero quelle che raggiungono \(q_j\) passando soltanto per stati da \(q_1\) a \(q_k\), ovvero ignorando lo stato \(q_{k+1}\), e quelle che attraversano almeno una volta \(q_{k+1}\).

  In particolare queste raggiungeranno \(q_{k+1}\) una prima volta, passando dunque solo per stati da \(q_1\) a \(q_k\), torneranno in \(q_{k+1}\) un numero arbitrario di volte, ancora passando solo per stati da \(q_1\) a \(q_k\), e infine lasceranno un'ultima volta \(q_{k+1}\) per raggiungere \(q_j\) passando ancora solo da stati tra \(q_1\) e \(q_k\). Questa divisione si può evincere più chiaramente dal diagramma seguente:

  \medskip
  \begin{center}
    \begin{tikzpicture}[initial text=, every state/.style={draw=blue!50,very thick,fill=blue!20}]
      \tikzset{dotsstyle/.style={draw=none,fill=none}}
      \node [state, initial, accepting]  (q_1) {\(q_i\)};
      \node [state, dotsstyle, right = of q_1] (q_4) {\(\cdots\cdots\)};
      \node [state, right = of q_4] (q_2) {\(q_{k+1}\)};
      \node [state, dotsstyle, below = 0.4 of q_2] (q_5) {\(\cdots\)};
      \node [state, dotsstyle, right = of q_2] (q_6) {\(\cdots\cdots \)};
      \node [state, dotsstyle, above = 0.2 of q_2] (q_7) {\(\cdots\cdots \)};
      \node [state, right = of q_6] (q_3) {\(q_j\)};

      % Text
      \node [below = -0.45 of q_4] {\(\cbrakets{q_1, \cdots, q_k}\)};
      \node [below = -0.45 of q_6] {\(\cbrakets{q_1, \cdots, q_k}\)};
      \node [above = -0.45 of q_7] {\(\cbrakets{q_1, \cdots, q_k}\)};
      \node [below = -0.45 of q_5] {\(\cbrakets{q_1, \cdots, q_k}\)};

      \path[-latex, every loop/.style = {-latex}]
      (q_4) edge [right] node [above] {} (q_2)
      (q_5) edge [out=0, in=310, below] node [below] {} (q_2)
      (q_6) edge [right] node [below] {} (q_3)
      (q_7) edge [out=0, in=135] node [below] {} (q_3);

      \path[]
      (q_1) edge [right] node [below] {} (q_4)
      (q_1) edge [out=45, in=180, left] node [below] {} (q_7)
      (q_2) edge [out=230, in=180, left] node [below] {} (q_5)
      (q_2) edge [right] node [above] {} (q_6);
    \end{tikzpicture}
  \end{center}
  \medskip

  Con queste tre osservazioni possiamo dimostriamo dunque per induzione su \(k\) che \textbf{ogni linguaggio } \(R^{(k)}_{i,j}\) può essere ottenuto da linguaggi finiti mediante un numero finito di operazioni di unione, concatenazione e iterazione.

  Possiamo ora facilmente vedere però che per ogni automa a stati finiti \(\mcM\) di \(n\) stati, in particolare quello considerato all'inizio della dimostrazione, vale che \(L(\mcM)= \bigcup\limits_{q_j \in F}R^{(n)}_{1,j}\).

  Ovvero il linguaggio accettato dall'automa è l'unione dei linguaggi che portano da \(1\) a un qualsiasi stato terminale fissato passando potenzialmente per tutti e \(n\) gli stati,
  cioè dunque un qualsiasi linguaggio regolare è sempre della forma \(R^{(k)}_{i,j}\) per qualche \(i,j,k\) definiti come sopra, la quale è sempre come da tesi.
\end{itemize}
\end{halfframedbox}

\newpage

\section{Espressioni Regolari}

Avendo analizzato le operazioni che permettono di ottenere linguaggi regolari, possiamo ora definire un modo per descrivere questi linguaggi in modo più compatto e semplice, ovvero delle stringhe che seguano specifiche regole sintattiche e che, abbinate a una semantica ben definita, permettano di descrivere tutti i linguaggi regolari.
\dfn{Espressioni regolari}{
  Dato \(A\) alfabeto finito, definisco:

  \begin{equation}
    \widehat{A} = \cbrakets{\underline{a} \vert a \in A} \cup \cbrakets{\underline{\cup}, \underline{\cdot}, \underline{*}, \underline{(}, \underline{)}, \underline{\varepsilon}, \underline{\varnothing}} % chktex 9
  \end{equation}

  Chiameremo \textbf{espressione regolare} su \(A\) una parola su \(\widehat{A}^*\) ottenuta dalle seguenti regole:

  \begin{enumerate}
    \item \(\forall a \in A: \underline{a}\) è un'espressione regolare, così come \(\underline{\varepsilon}\) e \(\underline{\varnothing}\).
    \item Se \(\alpha\) e \(\beta\) sono espressioni regolari, allora anche \(\underline{(}\alpha\underline{\cup}\beta\underline{)}\) lo è. % chktex 9
    \item Se \(\alpha\) e \(\beta\) sono espressioni regolari, allora anche \(\underline{(}\alpha\underline{\cdot}\beta\underline{)}\) lo è. % chktex 9
    \item Se \(\alpha\) è un'espressione regolare, allora anche \(\underline{(}\alpha\underline{)*}\) lo è. % chktex 9
    \item Non ci sono altre regole.
  \end{enumerate}
}

\subsection{Semantica delle Espressioni Regolari}

Oltre ad aver definito la sintassi delle espressioni regolari, per far si che queste possano essere utilizzate per descrivere linguaggi regolari è necessario definire una semantica per queste espressioni, ovvero un modo per interpretare le espressioni regolari come linguaggi regolari.

Se \(\gamma\) è un'espressione regolare su \(A\), il linguaggio associato a \(\gamma\): \(\abrakets{\gamma} \subseteq A^*\) è definito dalle regole semantiche:

\begin{enumerate}
  \item \(\forall a \in A: \abrakets{\underline{a}}=\cbrakets{a}\), inoltre \(\abrakets{\underline{\varepsilon}}=\cbrakets{\varepsilon}\) e \(\abrakets{\underline{\varnothing}}=\varnothing\).
  \item Se \(\gamma = \underline{(}\alpha\underline{\cup}\beta\underline{)}\), allora \(\abrakets{\gamma}=\abrakets{\alpha}\cup\abrakets{\beta}\). % chktex 9
  \item Se \(\gamma = \underline{(}\alpha\underline{\cdot}\beta\underline{)}\), allora \(\abrakets{\gamma}=\abrakets{\alpha}\cdot\abrakets{\beta}\).  % chktex 9
  \item Se \(\gamma = \underline{(}\alpha\underline{)}^{\underline{*}}\), allora \(\abrakets{\gamma}={(\abrakets{\alpha})}^*\). % chktex 9
\end{enumerate}

Il teorema di Kleene~\ref{Teorema di Kleene} può essere dunque reinterpretato come:

\begin{equation}
  L \subseteq A^* \text{ regolare } \iff \exists \alpha \text{ espressione regolare su } A \text{ tale che } L=\abrakets{\alpha}
\end{equation}

\begin{generalbox}
  [colframe=azure-gradient-3!90!black]
  {Esempio: Espressione Regolare}
  Consideriamo il linguaggio dell'esercizio~\ref{Automa L_a} sull'alfabeto \(\cbrakets{a,b}\).

  Si avrà che l'espressione regolare che lo definisce è:

  \begin{equation}
    \alpha = \rbrakets{\rbrakets{a \cup b } \cdot \rbrakets{\rbrakets{a \cup b} \cup  \rbrakets{a \cdot \rbrakets{\rbrakets{a \cup b}^* \cdot a}} \cup \rbrakets{b \cdot \rbrakets{\rbrakets{a\cup b}^* \cdot b}}}}
  \end{equation}
\end{generalbox}

Per comodità di scrittura, possiamo definire delle regole di precedenza tra gli operatori in modo da evitare di dover usare le parentesi e per sottintendere il simbolo \(\cdot\). In particolare, definiamo:

\begin{equation}
  * > \cdot > \cup
\end{equation}

Avremo quindi che l'espressione regolare dell'ultimo esempio può essere scritta come:

\begin{equation}
  \alpha = (a \cup b ) (a \cup b \cup a{(a\cup b)}^*a \cup b{(a\cup b)}^*b)
\end{equation}

\newpage

\section{Pumping Lemma per Linguaggi Regolari}

Fino ad ora abbiamo analizzato e descritto esclusivamente linguaggi regolari e come operare su di essi.

Vogliamo ora analizzare il problema opposto, ovvero come dimostrare che un linguaggio non è regolare.

Per fare ciò, introduciamo un risultato teorico che dimostra una condizione necessaria per la regolarità di un linguaggio, che se quindi negata ci permette di dimostrare che un linguaggio non è regolare.

Questo risultato è il \textbf{Pumping Lemma per Linguaggi Regolari}, e sfrutta un principio combinatorio detto \textbf{principio della piccionaia}.

Il \textbf{principio della piccionaia} afferma che se \(n\) piccioni vengono messi in \(m\) piccionaie, con \(n>m\), allora almeno una piccionaia conterrà almeno due piccioni.

Generalizzando questo concetto, potremmo dire che se ho \(n\) oggetti da distribuire su \(m\) insiemi, con \(n>m\), non potrò fare a meno di associare a più oggetti lo stesso insieme.

\begin{halfframedbox}{red!75!black}{Pumping lemma per linguaggi regolari (lemma uwu)}
  \textbf{Enunciato. } Sia \(\mcM\) un DFA con \(n\) stati, e \(x\in L(\mcM)\) tale che \(\abs{x} \geq n\). Allora esistono \(u,v,w \in A^*\) tali che:

  \begin{enumerate}
    \item \(x=uvw\)
    \item \(v \neq \varepsilon\)
    \item \(\forall i \in \mbN: uv^{i}w \in L(\mcM)\), ovvero \( \cbrakets{u}\cbrakets{v}^*\cbrakets{w} \subseteq L(\mcM)\)
  \end{enumerate}

  \begin{center}
    \textcolor{red!75!black}{\hdashrule[0.5ex]{18cm}{1pt}{3mm 3pt 1mm 2pt}}
  \end{center}

  \textbf{Dimostrazione. } Per il principio della piccionaia, nel percorso di accettazione di \(x\) (da \(q_1\) a \(q_k \in F\)), che passa necessariamente per almeno \(n+1\) stati dovendo includere anche lo stato iniziale, deve esistere un ciclo, ovvero si passerà per un certo stato \(q_l\) almeno due volte.

  Si ha allora \(x=uvw\), dove \(\delta^*(q_1,u)=q_l\), ovvero \(u\) è la parola che porta da \(q_1\) a \(q_l\), \(\delta^*(q_l,v)=q_l\), ovvero \(v\) è la parola descitta dal ciclo partente da \(q_l\), e \(\delta^*(q_l,w)=q_k\), ovvero \(w\) è la parola che porta dalla fine del ciclo fino allo stato terminale \(q_k\).

  Avremo dunque chiaramente che \(v \neq \varepsilon\), dovendo passare per \(q_l\) almeno due volte distinte, ma anche che \(\forall i \in \mbN: \delta^*(q_1,uv^{i}w)=q_k\), potendo percorrere il ciclo un numero arbitrario di volte, da cui la tesi.

  La suddivisione di \(x\) in \(u,v,w\) come da dimostrazione è più facilmente intuibile dal seguente diagramma:

  \medskip
  \begin{center}
    \begin{tikzpicture}[initial text=, every state/.style={draw=blue!50,very thick,fill=blue!20}]
      \tikzset{dotsstyle/.style={draw=none,fill=none}}

      \node [state, initial]  (q_1) {\(q_1\)};
      \node [state, dotsstyle, right = of q_1] (q_2) {\(\cdots\cdots\)};
      \node [state, right = of q_2] (q_3) {\(q_l\)};
      \node [state, dotsstyle, below = 0.4 of q_3] (q_4) {\(\cdots\)};
      \node [state, dotsstyle, right = of q_3] (q_5) {\(\cdots\cdots\)};
      \node [state, accepting, right = of q_5] (q_6) {\(q_k\)};

      % Text
      \node [below = -0.45 of q_2] {\(u\)};
      \node [below = -0.45 of q_4] {\(v\)};
      \node [below = -0.45 of q_5] {\(w\)};

      \path[-latex, every loop/.style = {-latex}]
      (q_1) edge [right] node [above] {} (q_2)
      (q_2) edge [right] node [above] {} (q_3)
      (q_4) edge [out=0, in=310, below] node [below] {} (q_3)
      (q_3) edge [out=230, in=180, left] node [below] {} (q_4)
      (q_3) edge [right] node [above] {} (q_5)
      (q_5) edge [right] node [above] {} (q_6);
    \end{tikzpicture}
  \end{center}

\end{halfframedbox}

Visto questo importante risultato, possiamo dunque dimostrare che un linguaggio non è regolare dimostrando che non rispetta il \textbf{Pumping Lemma}.

Oltre a fare ciò però è possibile derivare dal \textbf{Pumping Lemma} altri risultati teorici minori utili.

\begin{halfframedbox}{red!75!black}{Corollario 1}
  \textbf{Enunciato. } Sia \(\mcM\) un DFA con \(n\) stati, allora \(L(\mcM) \neq \varnothing \implies \exists x \in L(\mcM): \abs{x} < n\).

  \begin{center}
    \textcolor{red!75!black}{\hdashrule[0.5ex]{18cm}{1pt}{3mm 3pt 1mm 2pt}}
  \end{center}

  \textbf{Dimostrazione. } Per assurdo, supponiamo che \(x \in L(\mcM)\) sia la parola di lunghezza minima con \(\abs{x} \geq n\).

  Per il pumping lemma, esistono \(u,v,w \in A^*\) tali che \(x=uvw\), con \(v \neq \varepsilon\) e \(\forall i \in \mbN: uv^{i}w \in L(\mcM)\). Ma allora \(uv^0w=uw \in L(\mcM)\), contraddicendo la minimalità di \(x\).

\end{halfframedbox}

Questo corollario risulta particolarmente utile per determinare se, dati due DFA \(\mcM_1, \mcM_2\), si ha che \(L(\mcM_1) \subseteq L{(\mcM)}_2\).

Si ha infatti che, poiché \(L(\mcM_1) \setminus L(\mcM_2) = L(\mcM_1) \cap \overline{L(\mcM_2)}\)  è regolare, possiamo disegnare un DFA \(\mcM\) che accetti \(L(\mcM_1) \setminus L(\mcM_2)\), e verificare se \(L(\mcM)=\varnothing\).

Dal corollario appena dimostrato però, per dimostrare che la differenza sia vuota è sufficiente che
\begin{equation}
  \forall x \in A^* \abs{x} < n \implies x \nin L(\mcM)
\end{equation}

Ovvero che non ci siano parole di lunghezza minore del numero di stati dell'automa accettate dal linguaggio.

\newpage

\begin{halfframedbox}{red!75!black}{Corollario 2}
  \textbf{Enunciato. } Sia \(\mcM\) un DFA con \(n\) stati, allora \(L(\mcM)\) è infinito \(\iff \exists x \in L(\mcM): n \leq \abs{x} < 2n\).

  \begin{center}
    \textcolor{red!75!black}{\hdashrule[0.5ex]{18cm}{1pt}{3mm 3pt 1mm 2pt}}
  \end{center}

  \textbf{Dimostrazione. } Dimostriamo la doppia implicazione.

  \begin{itemize}
    \item [``\(\Rightarrow\)''] Dal \textbf{Pumping Lemma}, \(x=uvw\) con \(v \neq \varepsilon\) e \(\forall i \in \mbN: uv^{i}w \in L(\mcM)\). Avremo dunque che \(L\) accetta un numero infinito di parole, al variare di \(i\) in \(\mbN\).
    \item [``\(\Leftarrow\)''] Essendo \(L(\mcM)\) infinito, la lunghezza delle parole non può essere superiormente limitata, altrimenti sarebbe possibile contare il numero massimo di parole accettabili.

    Sia \(x\) dunque la più corta parola in \(L(\mcM)\) tale che \(\abs{x} \geq 2n\). Possiamo dunque scrivere \(x=yz\) con \(\abs{y}=n\) e \(\abs{z}\geq n\). Per il principio della piccionaia, avremo \(y=uvw\) con \(v \neq \varepsilon\) e \(\forall i \in \mbN: \delta^*(q_1,uv^{i}w)=\delta^*(q_1,y)\), ovvero al variare di \(i\) raggiungerò lo stesso stato raggiunto con \(y\), da cui \(\forall i \in \mbN: uv^{i}wz \in L(\mcM)\).

    Allora anche \(uwz \in L(\mcM)\) ma \(\abs{uwz} < \abs{yx}=\abs{x}\), da cui \(\abs{uwz} < 2n\) dalla minimalità di \(x\). Ma essendo \(\abs{z} \geq n\) per costruzione, si ha che \(\abs{uwz} \geq n\), da cui la tesi.
  \end{itemize}
\end{halfframedbox}

Vediamo ora finalmente degli esempi di \textbf{Linguaggi non Regolari}, e di come dimostrare ciò utilizzando il \textbf{Pumping Lemma}.

\begin{generalbox}[colframe=azure-gradient-3!90!black]{Esempi: Linguaggi non regolari}

  \begin{enumerate}
    \item Il linguaggio \(L = \cbrakets{a^{n}b^n \vert n \in \mbN}\), ovvero un linguaggio che accetti solo parole con un numero di \(a\) consecutive seguito da un ugual numero di \(b\) consecutive, \textbf{non} è regolare.

          \medskip

          Infatti sia per assurdo \(\mcM\) un DFA con \(p\) stati, tale che \(L(\mcM)=L\).

          \medskip

          Presa \(x=a^{p}b^p\), se \(x=uvw\) con \(v\neq \varepsilon\) sono possibili tre casi:
          \begin{enumerate}
            \item \(v=a^j, j>0 \implies uw=a^{p-j}b^p \nin L\)
            \item \(v=b^j, j>0 \implies uw=a^{p}b^{p-j} \nin L\)
            \item \(v=a^{j}b^k, j,k>0 \implies uv^2w = a^{p-j}{(a^{j}b^k)}^2b^{p-k}=a^{p}b^{k}a^{j}b^p \nin L\)
          \end{enumerate}

          \medskip

          In questo esempio specifico è possibile dimostrare la tesi anche con altri valori di \(i\), come ad esempio \(i=0\) nell'ultimo caso, ma in generale questo non è possibile.

          \medskip

          Abbiamo quindi che per qualsisi scelta di \(v\) il \textbf{Pumping Lemma} non è rispettato, dimostrando la tesi.

          \medskip

    \item Il linguaggio \(L' = \cbrakets{a^{m}b^n \vert m \geq n > 0}\), non è regolare.

          \medskip

          Infatti sia per assurdo \(\mcM\) un DFA con \(p\) stati, tale che \(L(\mcM)=L'\).

          \medskip

          Scelgo \(x=a^{p}b^{p} \in L, \abs{x}=2p\geq p\). Possiamo dividere i casi possibili in maniera analoga all'esempio precedente, con l'unica differenza che nel primo caso non potremo scegliere un'arbitraria \(i\) per la dimostrazione, poichè \( \forall i\geq 1\) la parola sarebbe ancora accettata dall'automa, dunque siamo costretti a scegliere \(i=0 \implies uv^{i}w = uw\) per dimostrare la tesi.

          \medskip

    \item Considerando invece il linguaggio \( L'' = \cbrakets{a^{m}b^n \vert m > n > 0}\), avremo che per la dimostrazione non potremo più usare la parola \(a^{p}b^{p} \nin L''\). Possiamo però considerare la parola \( x=a^{p}b^{p-1} \in L'', \abs{x}=2p-1\geq p\), che si suddivide nei medesimi tre casi degli esempi precedenti e si dimostra con la stessa scelta di \(i\) del secondo esempio.

          \medskip

    \item \textbf{Esercizio d'esame. } Mostrare che \(L = \cbrakets{a^{n}b^{n} \vert n \geq 0} \cup \abrakets{b^*\cdot a^*}\) non è regolare.

          \medskip

          \textbf{Soluzione. } Per quanto il \textbf{Pumping Lemma} sia uno strumento utile per dimostrare la non regolarità di un linguaggio, non è sempre l'approccio corretto da seguire. In questo caso andremo ad utilizzare le proprietà di chiusura dei linguaggi regolari.

          \medskip

          Per assurdo, sia \(L\) regolare. Allora dovrebbe esserlo anche \(L \cap \abrakets{a^*b^*}\), ma \(L \cap \abrakets{a^*b^*} = \cbrakets{a^{n}b^{n} \vert n \geq 0}\), che per quanto visto in precedenza non è regolare, da cui l'assurdo.
  \end{enumerate}

\end{generalbox}

Come già intuito dall'ultimo esempio, e come è evincibile dall'enunciato stesso, il \textbf{Pumping Lemma} non fornisce una condizione sufficiente per la regolarità di un linguaggio, bensì una condizione necessaria.

Possono esistere dunque linguaggi non regolari che rispettano il \textbf{Pumping Lemma}, e per questo motivo non è possibile dimostrare la regolarità di un linguaggio utilizzando il \textbf{Pumping Lemma}.

Utilizzando però l'enunciato di \textbf{Pumping Lemma} dato possiamo trovare dei linguaggi non regolari, per i quali anche la non regolarità non è dimostrabile utilizzandolo.

Vediamo dunque una versione più forte del \textbf{Pumping Lemma}, derivabile in maniera diretta dalla dimostrazione precedente, per poi vedere un esempio di linguaggio non regolare per il quale non è possibile dimostrare la non regolarità utilizzando il \textbf{Pumping Lemma} precedente.

\newpage

\begin{halfframedbox}{red!75!black}{Pumping Lemma per linguaggi regolari rafforzato}
  \textbf{Enunciato. } Sia \(\mcM\) un DFA con \(n\) stati, e \(x\in L(\mcM)\) tale che \(\abs{x} \geq n\). Allora esistono \(u,v,w \in A^*\) tali che:

  \begin{enumerate}
    \item \(x=uvw\)
    \item \(v \neq \varepsilon\)
    \item \(\forall i \in \mbN: uv^{i}w \in L\)
    \item \(\abs{uv} \leq n\)
    \end{enumerate}

  \begin{center}
    \textcolor{red!75!black}{\hdashrule[0.5ex]{18cm}{1pt}{3mm 3pt 1mm 2pt}}
  \end{center}

  \textbf{Dimostrazione. } Per dimostrare questa versione del \textbf{Pumping Lemma} consideriamo la parola \(x_n \subseteq x\) formata dalle prime \(n\) lettere.

  Per il principio della piccionaia sappiamo che già in questa parola deve essere presente almeno uno stato ripetuto.

  Potrò allora scegliere come \(u\) la sottoparola di \(x\) fino a questo primo stato ripetuto, con dunque \(u \subseteq x_n\), come \(v\) la parola del ciclo contenuto in \(x_n\), e come \(w\) la parola formata dalle restanti lettere di \(x\).

  Avendo scelto \(u,v \subseteq x_n\) si ha che \(\abs{uv} \leq n\), da cui la tesi.
\end{halfframedbox}

\begin{generalbox}[colframe=azure-gradient-3!90!black]{Esempio: Non regolarità del linguaggio dei palindromi}

  Vediamo dunque un esempio di linguaggio non regolare, che però non può essere dimostrato con il precedente \textbf{Pumping Lemma}.

  \medskip

  Sia \(L = \cbrakets{w \in \cbrakets{a,b}^* \vert w=\widetilde{w} } \), dove se \( w=w_1w_2\cdots w_n, w_1w_2\cdots w_n \in \cbrakets{a,b}^* \) allora \( \widetilde{w} = w_{n}w_{n-1}\cdots w_1 \). Questo linguaggio non è regolare.

  \medskip

  Proviamo inizialmente a dimostrare col \textbf{Pumping Lemma} originario che \(L\) non è regolare.

  \medskip

  Sia \(p\) il numero di stati di un DFA che accetta \(L\). Avremo che, per ogni parola di \(L\), si potrà sempre scegliere \(v\) come la sotto parola centrale e ottenere una parola in \(L\). Infatti, se \(x=uvw=\widetilde{x}\) con \(v=\widetilde{v}\), allora \(w=\widetilde{u}\) e dunque \(\forall i \in \mbN: uv^{i}w \in L\).

  \medskip

  Utilizzando il \textbf{Pumping Lemma} rafforzato però, sappiamo che \(\abs{uv} \leq p\), ovvero \(v\) dev'essere contenuta nelle prime \(p\) lettere di \(x\). Ma allora \(v\) non può essere la sotto parola centrale di \(x\).

  Consideriamo la parola \(x=a^{p}ba^{p} = uvw \), con \(v \neq \varepsilon \) e \(\abs{uv} \leq p\). Allora avremo necessariamente che:
  \begin{equation}
    v=a^k, p\geq k>0 \implies uw=a^{p-k}ba^{p} \nin L
  \end{equation}

  da cui la tesi.
\end{generalbox}
